\section{Performance tracking} \label{sec:tracking}

\noindent This chapter tracks the agreement of \ac{SMOS} \ac{L2} with three reference datasets -- \ac{ISMN} in situ \ac{FRMs}, ERA5-Land and \ac{SMAP} \ac{L3} (Chapter \ref{sec:datasets}) -- across the reports in this series, to detect unexpected changes in the \ac{SMOS} \ac{L2} data quality.
Two metrics are tracked: Pearson's R (only significant correlations, $p \leq$ \detokenize{$<ReportVars['Tracking R_between_0-ISMN_and_1-SMOS_L2']['tracking']['p_mask_thres']>$}) and the \ac{ubRMSD}.
Each report uses a 3-month window that ends with the month before the report date; consecutive monthly reports therefore share 2 months of data, and their results are not independent.
\\[1em]
\noindent\textbf{Expected range:} For each metric, the medians across all locations from up to \detokenize{$<ReportVars['Common']['n_epochs'] - 1>$} previous reports are used to fit a normal distribution.
The \detokenize{$<format(ReportVars['Tracking R_between_0-ISMN_and_1-SMOS_L2']['tracking']['alpha'] * 100, 'g')>$}\,\% and \detokenize{$<format(100 - ReportVars['Tracking R_between_0-ISMN_and_1-SMOS_L2']['tracking']['alpha'] * 100, 'g')>$}\,\% quantiles of this distribution define the expected range for the median of the current report.
A median outside this range is flagged (\textcolor{red}{red}).
At least \detokenize{$<ReportVars['Tracking R_between_0-ISMN_and_1-SMOS_L2']['tracking']['min_history']>$} previous reports are required to derive an expected range; otherwise the metric is not classified (\textcolor{gray}{gray}).
\\[1em]
\noindent\textbf{Change at common locations:} The median across all locations also depends on \textit{which} locations are included, which can change between reports.
Therefore, the median change of the metric at the locations shared with the previous report is given in addition ($\Delta$ prev.).
It is informative only and does not affect the tracking status.

\begin{table}[ht!]
\centering
\small
\caption{Tracking results for this report. Median: median across all valid locations. $\Delta$ prev.: median change at the locations shared with the previous report. Status: \textcolor{green!60!black}{green}: within the expected range, \textcolor{red}{red}: outside, \textcolor{gray}{gray}: not classified.}
\label{tab:tracking_status}
\begin{tabular}{@{}llrcrc@{}}
\toprule
Reference & Metric & Median & Expected range & $\Delta$ prev. & Status \\
\midrule
\ac{ISMN} \ac{FRMs} & R & \detokenize{$<'--' if ReportVars['Tracking R_between_0-ISMN_and_1-SMOS_L2']['tracking']['current_median'] is None else format(ReportVars['Tracking R_between_0-ISMN_and_1-SMOS_L2']['tracking']['current_median'], '.2f')>$} & \detokenize{$<'--' if ReportVars['Tracking R_between_0-ISMN_and_1-SMOS_L2']['tracking']['expected_lower'] is None else f"[{ReportVars['Tracking R_between_0-ISMN_and_1-SMOS_L2']['tracking']['expected_lower']:.2f}, {ReportVars['Tracking R_between_0-ISMN_and_1-SMOS_L2']['tracking']['expected_upper']:.2f}]">$} & \detokenize{$<'--' if ReportVars['Tracking R_between_0-ISMN_and_1-SMOS_L2']['tracking']['paired_previous']['median_diff'] is None else format(ReportVars['Tracking R_between_0-ISMN_and_1-SMOS_L2']['tracking']['paired_previous']['median_diff'], '+.3f')>$} & \circlemark{\detokenize{$<ReportVars['Tracking R_between_0-ISMN_and_1-SMOS_L2']['tracking']['tracking_color']>$}} \\
 & \ac{ubRMSD} & \detokenize{$<'--' if ReportVars['Tracking urmsd_between_0-ISMN_and_1-SMOS_L2']['tracking']['current_median'] is None else format(ReportVars['Tracking urmsd_between_0-ISMN_and_1-SMOS_L2']['tracking']['current_median'], '.3f')>$} & \detokenize{$<'--' if ReportVars['Tracking urmsd_between_0-ISMN_and_1-SMOS_L2']['tracking']['expected_lower'] is None else f"[{ReportVars['Tracking urmsd_between_0-ISMN_and_1-SMOS_L2']['tracking']['expected_lower']:.3f}, {ReportVars['Tracking urmsd_between_0-ISMN_and_1-SMOS_L2']['tracking']['expected_upper']:.3f}]">$} & \detokenize{$<'--' if ReportVars['Tracking urmsd_between_0-ISMN_and_1-SMOS_L2']['tracking']['paired_previous']['median_diff'] is None else format(ReportVars['Tracking urmsd_between_0-ISMN_and_1-SMOS_L2']['tracking']['paired_previous']['median_diff'], '+.3f')>$} & \circlemark{\detokenize{$<ReportVars['Tracking urmsd_between_0-ISMN_and_1-SMOS_L2']['tracking']['tracking_color']>$}} \\
\midrule
ERA5-Land & R & \detokenize{$<'--' if ReportVars['Tracking R_between_0-SMOS_L2_and_1-ERA5_LAND']['tracking']['current_median'] is None else format(ReportVars['Tracking R_between_0-SMOS_L2_and_1-ERA5_LAND']['tracking']['current_median'], '.2f')>$} & \detokenize{$<'--' if ReportVars['Tracking R_between_0-SMOS_L2_and_1-ERA5_LAND']['tracking']['expected_lower'] is None else f"[{ReportVars['Tracking R_between_0-SMOS_L2_and_1-ERA5_LAND']['tracking']['expected_lower']:.2f}, {ReportVars['Tracking R_between_0-SMOS_L2_and_1-ERA5_LAND']['tracking']['expected_upper']:.2f}]">$} & \detokenize{$<'--' if ReportVars['Tracking R_between_0-SMOS_L2_and_1-ERA5_LAND']['tracking']['paired_previous']['median_diff'] is None else format(ReportVars['Tracking R_between_0-SMOS_L2_and_1-ERA5_LAND']['tracking']['paired_previous']['median_diff'], '+.3f')>$} & \circlemark{\detokenize{$<ReportVars['Tracking R_between_0-SMOS_L2_and_1-ERA5_LAND']['tracking']['tracking_color']>$}} \\
 & \ac{ubRMSD} & \detokenize{$<'--' if ReportVars['Tracking urmsd_between_0-SMOS_L2_and_1-ERA5_LAND']['tracking']['current_median'] is None else format(ReportVars['Tracking urmsd_between_0-SMOS_L2_and_1-ERA5_LAND']['tracking']['current_median'], '.3f')>$} & \detokenize{$<'--' if ReportVars['Tracking urmsd_between_0-SMOS_L2_and_1-ERA5_LAND']['tracking']['expected_lower'] is None else f"[{ReportVars['Tracking urmsd_between_0-SMOS_L2_and_1-ERA5_LAND']['tracking']['expected_lower']:.3f}, {ReportVars['Tracking urmsd_between_0-SMOS_L2_and_1-ERA5_LAND']['tracking']['expected_upper']:.3f}]">$} & \detokenize{$<'--' if ReportVars['Tracking urmsd_between_0-SMOS_L2_and_1-ERA5_LAND']['tracking']['paired_previous']['median_diff'] is None else format(ReportVars['Tracking urmsd_between_0-SMOS_L2_and_1-ERA5_LAND']['tracking']['paired_previous']['median_diff'], '+.3f')>$} & \circlemark{\detokenize{$<ReportVars['Tracking urmsd_between_0-SMOS_L2_and_1-ERA5_LAND']['tracking']['tracking_color']>$}} \\
\midrule
\ac{SMAP} \ac{L3} & R & \detokenize{$<'--' if ReportVars['Tracking R_between_0-SMOS_L2_and_1-SMAP_L3']['tracking']['current_median'] is None else format(ReportVars['Tracking R_between_0-SMOS_L2_and_1-SMAP_L3']['tracking']['current_median'], '.2f')>$} & \detokenize{$<'--' if ReportVars['Tracking R_between_0-SMOS_L2_and_1-SMAP_L3']['tracking']['expected_lower'] is None else f"[{ReportVars['Tracking R_between_0-SMOS_L2_and_1-SMAP_L3']['tracking']['expected_lower']:.2f}, {ReportVars['Tracking R_between_0-SMOS_L2_and_1-SMAP_L3']['tracking']['expected_upper']:.2f}]">$} & \detokenize{$<'--' if ReportVars['Tracking R_between_0-SMOS_L2_and_1-SMAP_L3']['tracking']['paired_previous']['median_diff'] is None else format(ReportVars['Tracking R_between_0-SMOS_L2_and_1-SMAP_L3']['tracking']['paired_previous']['median_diff'], '+.3f')>$} & \circlemark{\detokenize{$<ReportVars['Tracking R_between_0-SMOS_L2_and_1-SMAP_L3']['tracking']['tracking_color']>$}} \\
 & \ac{ubRMSD} & \detokenize{$<'--' if ReportVars['Tracking urmsd_between_0-SMOS_L2_and_1-SMAP_L3']['tracking']['current_median'] is None else format(ReportVars['Tracking urmsd_between_0-SMOS_L2_and_1-SMAP_L3']['tracking']['current_median'], '.3f')>$} & \detokenize{$<'--' if ReportVars['Tracking urmsd_between_0-SMOS_L2_and_1-SMAP_L3']['tracking']['expected_lower'] is None else f"[{ReportVars['Tracking urmsd_between_0-SMOS_L2_and_1-SMAP_L3']['tracking']['expected_lower']:.3f}, {ReportVars['Tracking urmsd_between_0-SMOS_L2_and_1-SMAP_L3']['tracking']['expected_upper']:.3f}]">$} & \detokenize{$<'--' if ReportVars['Tracking urmsd_between_0-SMOS_L2_and_1-SMAP_L3']['tracking']['paired_previous']['median_diff'] is None else format(ReportVars['Tracking urmsd_between_0-SMOS_L2_and_1-SMAP_L3']['tracking']['paired_previous']['median_diff'], '+.3f')>$} & \circlemark{\detokenize{$<ReportVars['Tracking urmsd_between_0-SMOS_L2_and_1-SMAP_L3']['tracking']['tracking_color']>$}} \\
\bottomrule
\end{tabular}
\end{table}

\noindent Figs. \ref{fig:tracking_ismn} to \ref{fig:tracking_smap} show the distribution of the metrics across all locations for each report (box: 25th to 75th percentile, whiskers: 5th to 95th percentile, black line and dots: median, n: number of locations).
The current report is the rightmost box, colored according to its tracking status; the shaded area is its expected median range.
All values are also listed in Chapter \ref{sec:tracking_tables}.

\FloatBarrier

\subsection{\texorpdfstring{\acs{ISMN} \acs{FRMs}}{ISMN FRMs}} \label{ssec:tracking_ismn}

\noindent The number of stations with valid data varies between reports (e.g. frozen soils in winter, network data gaps); see Table \ref{tab:networks} for the sample of this report.

\begin{itemize}
    \item \textbf{R:} \detokenize{$<ReportVars['Tracking R_between_0-ISMN_and_1-SMOS_L2']['tracking']['summary']>$}
    \item \textbf{\ac{ubRMSD}:} \detokenize{$<ReportVars['Tracking urmsd_between_0-ISMN_and_1-SMOS_L2']['tracking']['summary']>$}
\end{itemize}

\begin{figure}[ht!]
    \centering
    \begin{subfigure}{0.8\textwidth}
        \centering
        \includegraphics[width=\textwidth]{./tracking/tracking_R_between_0-ISMN_and_1-SMOS_L2.png}
        \caption{Pearson's R}
    \end{subfigure}
    \begin{subfigure}{0.8\textwidth}
        \centering
        \includegraphics[width=\textwidth]{./tracking/tracking_urmsd_between_0-ISMN_and_1-SMOS_L2.png}
        \caption{\ac{ubRMSD}}
    \end{subfigure}
    \caption{Cross-report tracking of Pearson's R (a) and \ac{ubRMSD} (b) between \ac{SMOS} \ac{L2} and \ac{ISMN} \ac{FRMs}.}
    \label{fig:tracking_ismn}
\end{figure}

\FloatBarrier

\subsection{\texorpdfstring{ERA5-Land}{ERA5-Land}} \label{ssec:tracking_e5l}

\noindent The spatially complete ERA5-Land reanalysis provides the largest and most stable sample for tracking.

\begin{itemize}
    \item \textbf{R:} \detokenize{$<ReportVars['Tracking R_between_0-SMOS_L2_and_1-ERA5_LAND']['tracking']['summary']>$}
    \item \textbf{\ac{ubRMSD}:} \detokenize{$<ReportVars['Tracking urmsd_between_0-SMOS_L2_and_1-ERA5_LAND']['tracking']['summary']>$}
\end{itemize}

\begin{figure}[ht!]
    \centering
    \begin{subfigure}{0.8\textwidth}
        \centering
        \includegraphics[width=\textwidth]{./tracking/tracking_R_between_0-SMOS_L2_and_1-ERA5_LAND.png}
        \caption{Pearson's R}
    \end{subfigure}
    \begin{subfigure}{0.8\textwidth}
        \centering
        \includegraphics[width=\textwidth]{./tracking/tracking_urmsd_between_0-SMOS_L2_and_1-ERA5_LAND.png}
        \caption{\ac{ubRMSD}}
    \end{subfigure}
    \caption{Cross-report tracking of Pearson's R (a) and \ac{ubRMSD} (b) between \ac{SMOS} \ac{L2} and ERA5-Land.}
    \label{fig:tracking_e5l}
\end{figure}

\FloatBarrier

\subsection{\texorpdfstring{\acs{SMAP} \acs{L3}}{SMAP L3}} \label{ssec:tracking_smap}

\noindent A change that is only seen against \ac{SMAP} can point to issues in either of the two satellite products.

\begin{itemize}
    \item \textbf{R:} \detokenize{$<ReportVars['Tracking R_between_0-SMOS_L2_and_1-SMAP_L3']['tracking']['summary']>$}
    \item \textbf{\ac{ubRMSD}:} \detokenize{$<ReportVars['Tracking urmsd_between_0-SMOS_L2_and_1-SMAP_L3']['tracking']['summary']>$}
\end{itemize}

\begin{figure}[ht!]
    \centering
    \begin{subfigure}{0.8\textwidth}
        \centering
        \includegraphics[width=\textwidth]{./tracking/tracking_R_between_0-SMOS_L2_and_1-SMAP_L3.png}
        \caption{Pearson's R}
    \end{subfigure}
    \begin{subfigure}{0.8\textwidth}
        \centering
        \includegraphics[width=\textwidth]{./tracking/tracking_urmsd_between_0-SMOS_L2_and_1-SMAP_L3.png}
        \caption{\ac{ubRMSD}}
    \end{subfigure}
    \caption{Cross-report tracking of Pearson's R (a) and \ac{ubRMSD} (b) between \ac{SMOS} \ac{L2} and \ac{SMAP} \ac{L3}.}
    \label{fig:tracking_smap}
\end{figure}

\FloatBarrier

\subsection{Data availability} \label{ssec:tracking_coverage}

\noindent Two indicators show whether the amount of data the metrics are based on changes between reports.
They are informative only and do not affect the tracking status.
\\[1em]
\noindent\textbf{Spatial coverage:} Fig. \ref{fig:coverage} shows the number of grid points for which the validation was successful (status GOOD) against the two gridded references, ERA5-Land and \ac{SMAP} \ac{L3}, for each report.
Changes in the coverage can indicate changes in the data availability of \ac{SMOS} \ac{L2} or the reference (e.g. due to frozen soils in winter or \ac{RFI}), and affect how comparable the metrics of different reports are.
The number of valid \ac{ISMN} sensors is given in Table \ref{tab:networks}.
\detokenize{$<ReportVars['Coverage']['coverage']['summary']>$}
\\[1em]
\noindent\textbf{Number of observations:} Fig. \ref{fig:tracking_nobs} shows the distribution of the number of temporally matched observations (n\_obs) per grid point in the validation against ERA5-Land.
As ERA5-Land is gap-free, n\_obs reflects the availability of \ac{SMOS} \ac{L2} (the temporal reference).
A drop in n\_obs at an unchanged spatial coverage indicates a temporal gap in the \ac{SMOS} \ac{L2} data stream (e.g. an outage that produced empty files): the metrics are then still computed at the same locations, but from fewer observations, and are less robust.
\detokenize{$<ReportVars['Nobs']['tracking']['summary']>$}

\begin{figure}[ht!]
    \centering
    \includegraphics[width=0.85\textwidth]{./tracking/coverage.png}
    \caption{Number of grid points with successful validation (status GOOD) per report against ERA5-Land and \ac{SMAP} \ac{L3}. The numbers are also listed in Table \ref{tab:stats_coverage}.}
    \label{fig:coverage}
\end{figure}

\begin{figure}[ht!]
    \centering
    \includegraphics[width=0.85\textwidth]{./tracking/tracking_n_obs.png}
    \caption{Cross-report tracking of the number of temporally matched observations (n\_obs) per grid point between \ac{SMOS} \ac{L2} and ERA5-Land. Shaded area: expected range derived from the previous reports.}
    \label{fig:tracking_nobs}
\end{figure}

\FloatBarrier

